% !TEX root = task-1-15-6.tex
% Editable task bank. Append new BankTask blocks at the end of the document.
\documentclass[12pt,a4paper]{article}
\usepackage{physics-textbook}
\usepackage{graphicx,needspace}
\newif\ifShowSolutions
\ifdefined\TutorEdition
  \ShowSolutionstrue
\else
  \ShowSolutionsfalse % Change to \ShowSolutionstrue to display solutions.
\fi
\newcounter{banktask}
\newcommand{\BankTask}[3]{%
  \refstepcounter{banktask}%
  \Needspace{35mm}\par\medskip\noindent\begin{minipage}{\linewidth}%
  \noindent\textbf{\thebanktask.}\quad #2\par
  \ifShowSolutions
    #3\par
    \ifdefined\TutorEdition
      \ifcsname TutorSource#1\endcsname
        \smallskip{\small\csname TutorSource#1\endcsname\par}%
      \fi
    \fi
  \fi
  \end{minipage}\par\medskip
}
\begin{document}
\BookSubsection{1.15.6}{Потенциальная энергия пружины}


% Source sdamgia-415: https://phys-ege.sdamgia.ru/problem?id=415
\BankTask{sdamgia-415}{
Растянутая на $2\,\text{см}$ стальная пружина обладает потенциальной энергией упругой деформации $4\,\text{Дж}$. На сколько увеличится потенциальная энергия упругой деформации при растяжении этой пружины еще на $2\,\text{см}$?
}{
\smallskip\noindent\textbf{Решение}\par
Потенциальная энергия упругой деформации пропорциональна квадрату удлинения пружины: $E(x)=\frac{kx^2}{2}=4\,\text{Дж}.$ Увеличение деформации в $2$ раза приведет к увеличению энергии в $4$ раза: $E(2x)=\frac{k(2x)^2}{2}=4\cdot\frac{kx^2}{2}=16\,\text{Дж},$ то есть потенциальная энергия пружины увеличится на $16-4=12\,\text{Дж}.$
 
Ответ: $12$.
}

% Source sdamgia-537: https://phys-ege.sdamgia.ru/problem?id=537
\BankTask{sdamgia-537}{
Первая пружина имеет жесткость $20\,\text{Н/м},$ вторая — $40\,\text{Н/м}.$ Обе пружины растянуты на $1\,\text{см}$. Чему равно отношение потенциальных энергии пружин $\frac{E_2}{E_1}$?
}{
\smallskip\noindent\textbf{Решение}\par
Потенциальная энергия деформированной пружины пропорциональна ее жесткости. Следовательно, отношение потенциальных энергии пружин равно:


$\frac{E_2}{E_1}=\frac{k_2}{k_1}=\frac{40\,\text{Н/м}}{20\,\text{Н/м}}=2.$

Ответ: $2$.
}

% Source sdamgia-41158: https://phys-ege.sdamgia.ru/problem?id=41158
\BankTask{sdamgia-41158}{
При удлинении пружины на $4\,\text{см}$ из недеформированного состояния её потенциальная энергия становится равной $0{,}6\,\text{Дж}$. Чему будет равна потенциальная энергия этой пружины, если из недеформированного состояния сжать её на $2\,\text{см}$?
}{
\smallskip\noindent\textbf{Решение}\par
В первом случае потенциальная энергия упруго деформированной пружины равна $E_{\text{п}1}=\frac{kx_1^2}{2},$ во втором — $E_{\text{п}2}=\frac{0{,}25kx_1^2}{2},$ откуда находим:


$E_{\text{п}2}=0{,}25E_{\text{п}1}=0{,}25\cdot0{,}6=0{,}15\,\text{Дж}.$

Ответ: $0{,}15$.
}

% Source sdamgia-29977: https://phys-ege.sdamgia.ru/problem?id=29977
\BankTask{sdamgia-29977}{
Потенциальная энергия упругой пружины при ее растяжении на $2\,\text{см}$ равна $2\,\text{Дж}$. Найдите модуль изменения потенциальной энергии этой пружины при уменьшении ее растяжения на $0{,}5\,\text{см}$.
}{
\smallskip\noindent\textbf{Решение}\par
Потенциальная энергия упруго деформированной пружины $E_p=\frac{kx^2}{2},$ модуль изменения потенциальной энергии $|E_p|=\frac k2(x_1^2-x_2^2).$ Учитывая, что удлинение пружины во втором случае равно $1{,}5\,\text{см}$, получаем:


$|E_p|=E_{p1}\left(1-\left(\frac{x_2}{x_1}\right)^2\right)=2\left(1-\left(\frac{1{,}5}{2}\right)^2\right)=0{,}875\,\text{Дж}.$ 

Ответ: $0{,}875$.
}

% Source sdamgia-31877: https://phys-ege.sdamgia.ru/problem?id=31877
\BankTask{sdamgia-31877}{
Жесткость пружины детского игрушечного пистолета равна $40\,\text{Н/м}$. Для его зарядки нужно приложить максимальную силу $2\,\text{Н}$, направленную вдоль оси пружины. Какую кинетическую энергию приобретает шарик, которым выстреливают из этого пистолета, если трение в механизме пистолета очень мало?
}{
\smallskip\noindent\textbf{Решение}\par
При выстреле происходит превращение потенциальной энергии упруго деформированной пружины $E_{\text{п}}=\frac{kx^2}{2}$ в кинетическую энергию шарика. По закону сохранения энергии $E_{\text{п}}=E_{\text{к}}.$
Деформация пружины из закона Гука $x=\frac Fk,$ откуда получаем:


$E_{\text{к}}=\frac{F^2}{2k}=\frac4{2\cdot40}=0{,}05\,\text{Дж}=50\,\text{мДж}.$

Ответ: $50$.
}

% Source sdamgia-30484: https://phys-ege.sdamgia.ru/problem?id=30484
\BankTask{sdamgia-30484}{
Груз массой $0{,}5\,\text{кг}$ подвешен к концу легкой пружины жесткостью $100\,\text{Н/м}$, второй конец которой прикреплен к потолку. Сначала этот груз удерживают так, что пружина не деформирована и расположена вертикально. Затем груз отпускают без начальной скорости. Найдите максимальное удлинение пружины после начала движения груза.
}{
\smallskip\noindent\textbf{Решение}\par
Примем за нулевой уровень высоты точку, в которой груз остановится. Тогда в верхней точке пружина не деформирована, груз обладает потенциальной энергией $E_{p1}=mgh.$ В нижней точке, в которой груз остановится, пружина будет обладать потенциальной энергией $E_{p2}=\frac{kh^2}{2}.$ Тогда по закону сохранения энергии $mgh=\frac{kh^2}{2},$ откуда: 


$h=\frac{2mg}{k}=\frac{2\cdot0{,}5\cdot10}{100}=0{,}1\,\text{м}=10\,\text{см}.$

Ответ: $10$.
}

% Source sdamgia-51009: https://phys-ege.sdamgia.ru/problem?id=51009
\BankTask{sdamgia-51009}{
Ученик исследовал зависимость силы упругости  $F$ от её растяжения $x$ и получил следующие результаты:
 
\begin{center}\begin{tabular}{|c|c|c|c|c|c|c|}\hline
$F$, $\text{Н}$&$0$&$0{,}5$&$1$&$1{,}5$&$2$&$2{,}5$\\\hline
$x$, $\text{м}$&$0$&$0{,}02$&$0{,}04$&$0{,}06$&$0{,}08$&$0{,}10$\\\hline
\end{tabular}\end{center}
Определите потенциальную энергию пружины при её растяжении на $0{,}04\,\text{м}$.
}{
\smallskip\noindent\textbf{Решение}\par
Потенциальная энергия пружины $E_{\text{п}}=\frac{kx^2}{2},$ при этом жесткость пружины $k=\frac Fx.$ Тогда находим:


$E_{\text{п}}=\frac{Fx}{2}=\frac{1\cdot0{,}04}{2}=0{,}02\,\text{Дж}.$

Ответ: $0{,}02$
}

% Source sdamgia-29461: https://phys-ege.sdamgia.ru/problem?id=29461
\BankTask{sdamgia-29461}{
Нить, удерживающая легкую пружину в сжатом состоянии, внезапно оборвалась (см. рисунок). На сколько сантиметров была сжата пружина жесткостью $20\,\text{кН/м}$, если она сообщила бруску массой $20\,\text{г}$ максимальную скорость $10\,\text{м/с}$? Трением бруска о поверхность и массой пружины пренебречь.
\begin{center}\begin{tikzpicture}[task graph]
\draw[task curve] (-.8,0)--(4.6,0)--(4.6,1.85);
\foreach \x in {-.7,-.45,...,4.5}{\draw[task tick] (\x,0)--++(.2,-.2);}
\foreach \y in {.1,.35,...,1.7}{\draw[task tick] (4.6,\y)--++(.2,-.2);}
\draw[task curve,fill=orange!25] (-.7,0) rectangle (.5,1.2);
\draw[task curve] (.5,.6)--(.65,.6)--(.8,1.2)--(1.1,0)--(1.4,1.2)--(1.7,0)--(2,1.2)--(2.3,0)--(2.6,1.2)--(2.9,0)--(3.2,1.2)--(3.5,0)--(3.8,1.2)--(4.1,0)--(4.25,.6)--(4.6,.6);
\draw[task curve,gray!65] (.5,.6)--(.65,.6)--(.8,0)--(1.1,1.2)--(1.4,0)--(1.7,1.2)--(2,0)--(2.3,1.2)--(2.6,0)--(2.9,1.2)--(3.2,0)--(3.5,1.2)--(3.8,0)--(4.1,1.2)--(4.25,.6)--(4.6,.6);
\draw[task tick,gray!65] (.5,.6)--(4.6,.6);
\end{tikzpicture}\end{center}
}{
\smallskip\noindent\textbf{Решение}\par
В начальном положении в системе пружина обладала потенциальной энергией $E_p=\frac{kx^2}{2},$ брусок находился в покое. Далее брусок приобретет максимальную скорость и кинетическую энергию $E_k=\frac{mv^2}{2},$ когда пружина будет не деформирована. Учитывая, что работа сил тяжести и реакции опоры равна $0$, сила трения отсутствует, применяем закон сохранения энергии $E_p=E_k$ или $\frac{kx^2}{2}=\frac{mv^2}{2}.$ Откуда первоначальная деформация пружины: 


$x=v\sqrt{\frac mk}=10\sqrt{\frac{20\cdot10^{-3}}{20\cdot10^3}}=10^{-2}\,\text{м}=1\,\text{см}.$

Ответ: $1\,\text{см}$.
}

% Source sdamgia-4745: https://phys-ege.sdamgia.ru/problem?id=4745
\BankTask{sdamgia-4745}{
Горизонтально расположенная невесомая пружина с жесткостью $k=1000\,\text{Н/м}$ находится в недеформированном состоянии. Один ее конец закреплен, а другой касается бруска массой $M=0{,}1\,\text{кг},$ находящегося на горизонтальной поверхности. Брусок сдвигают, сжимая пружину, и отпускают. На какую длину $\Delta x$ была сжата пружина, если после отпускания бруска его скорость достигла величины $v=1\,\text{м/с}$? Трение не учитывать.
}{
\smallskip\noindent\textbf{Решение}\par
После того, как пружину сжали, совершив над ней работу, на систему «пружина — брусок» больше не действует внешних сил, совершающих работу, а значит, для этой системы выполняется закон сохранения полной механической энергии $E_{\text{пот}}=E_{\text{кин}},$ при этом потенциальная энергия упруго деформированной пружины равна $E_{\text{пот}}=\frac{k\Delta x^2}{2},$ кинетическая энергия бруска $E_{\text{кин}}=\frac{Mv^2}{2}.$ Таким образом, получаем:


$\frac{Mv^2}{2}=\frac{k(\Delta x)^2}{2}$
$\Delta x=v\sqrt{\frac Mk}=1\,\text{м/с}\cdot\sqrt{\frac{0{,}1\,\text{кг}}{1000\,\text{Н/м}}}=0{,}01\,\text{м}.$

Ответ: $0{,}01\,\text{м}$.
}

% Source sdamgia-10234: https://phys-ege.sdamgia.ru/problem?id=10234
\BankTask{sdamgia-10234}{
Упругая легкая пружина жесткостью $80\,\text{Н/м}$ одним концом прикреплена к лапке штатива. К свободному концу пружины подвешен груз массой $200\,\text{г}$. Определите потенциальную энергию растянутой пружины.
}{
\smallskip\noindent\textbf{Решение}\par
Под действием силы тяжести пружина будет растягиваться до тех пор, пока сила упругости, возникшая при деформации пружины, равная $F_{\text{у}}=k\Delta x,$ не уравновесит силу тяжести $F_{\text{у}}=mg,$ откуда деформация пружины $\Delta x=\frac{mg}{k}.$ Потенциальная энергия пружины может быть найдена по формуле:


$E=\frac{k\Delta x^2}{2}=\frac{(mg)^2}{2k}=\frac{(0{,}2\cdot10)^2}{2\cdot80}=25\,\text{мДж}.$

Ответ: $25$ мДж.
}


% Source ADA046: https://ege.fipi.ru/bank/questions.php?search=1&proj=BA1F39653304A5B041B656915DC36B38&qid=ADA046
\BankTask{ADA046}{
При упругой деформации, равной $3\,\text{см}$, потенциальная энергия пружины равна $6\,\text{Дж}$. На сколько изменится потенциальная энергия этой пружины при увеличении упругой деформации ещё на $3\,\text{см}$? Выразите ответ в джоулях.
}{
\smallskip\noindent\textbf{Решение от GPT}\par
$E_{\text{п}}=kx^2/2$. Деформация увеличилась вдвое, поэтому $E_2=4E_1=24\,\text{Дж}$. Изменение энергии $\Delta E=E_2-E_1=24-6=18\,\text{Дж}$. Ответ: увеличится на $18\,\text{Дж}$.
}

% Source 1C4DA4: https://ege.fipi.ru/bank/questions.php?search=1&proj=BA1F39653304A5B041B656915DC36B38&qid=1C4DA4
\BankTask{1C4DA4}{
Первая пружина имеет жёсткость $20\,\text{Н}/\text{м}$, вторая — $40\,\text{Н}/\text{м}$. Первая растянута на $2\,\text{см}$, вторая — на $1\,\text{см}$. Чему равно отношение потенциальной энергии второй пружины к потенциальной энергии первой $\frac{E_2}{E_1}?$
}{
\smallskip\noindent\textbf{Решение от GPT}\par
$E_2/E_1=k_2x_2^2/(k_1x_1^2)=(40/20)(1/2)^2=0{,}5$.
}

% Source 1F9FAE: https://ege.fipi.ru/bank/questions.php?search=1&proj=BA1F39653304A5B041B656915DC36B38&qid=1F9FAE
\BankTask{1F9FAE}{
Жёсткость пружины равна $7500\,\text{Н}/\text{м}$. Какова энергия упругой деформации этой пружины при её растяжении на $4\,\text{см}$? Результат выразите в $\text{Дж}$.
}{
\smallskip\noindent\textbf{Решение от GPT}\par
$x=0{,}04\,\text{м}$. Энергия $E=kx^2/2=7500\cdot0{,}04^2/2=6\,\text{Дж}$.
}

% Source 38979: https://phys-ege.sdamgia.ru/problem?id=38979
\BankTask{38979}{
\begin{center}\begin{tikzpicture}[task graph,>=stealth]\draw(-.3,0)--(6,0)--(6,2);\foreach \x in {-.3,-.1,...,5.9}{\draw(\x,0)--++(.15,-.15);}\foreach \y in {.1,.3,...,1.9}{\draw(6,\y)--++(.15,.15);}\draw[fill=gray!15](0,0)rectangle(.6,.58)node[midway]{\(m\)};\draw[fill=gray!15](2,0)rectangle(3.15,1.15)node[midway]{\(M\)};\draw(4.85,0)--(4.85,.9);\draw(4.85,.68)--(4.95,.82)--(5.12,.57)--(5.29,.82)--(5.46,.57)--(5.63,.82)--(5.80,.57)--(6,.82);\node[above]at(5.45,.85){\(k\)};\draw(3.15,0)--(3.15,-.85)(4.85,0)--(4.85,-.85);\draw[<->](3.15,-.58)--(4.85,-.58)node[midway,below]{\(L\)};\draw[->] (.6,.31)--(1.45,.31) node[above]{\(\vec v\)};\end{tikzpicture}\end{center} Небольшое тело массой m  =  $100\,\text{г}$, скользящее по гладкой горизонтальной поверхности, абсолютно неупруго сталкивается с неподвижным телом массой M  =  $2$m. При дальнейшем поступательном движении тела налетают на недеформированную пружину, одним концом прикреплённую к стене (см. рис.). Через какое время t после абсолютно неупругого удара тела вернутся в точку столкновения? Скорость движения тела до столкновения $v = 3 \,\text{м/с}$ жёсткость пружины k  =  $30\,\text{Н/м}$, а расстояние от точки столкновения до пружины L  =  $15\,\text{см}$.
}{
\smallskip\noindent\textbf{Решение}\par
Обоснование.
Рассмотрим решение задачи в системе отсчета, связанной с Землей, которую можно считать инерциальной. Тела движутся поступательно, поэтому их движение можно описывать моделью материальной точки.
При отсутствии силы трения система тел является замкнутой, поэтому в ИСО можем применить закон сохранения импульса тел.
Так как поверхность гладкая, работа внешних сил при движении по горизонтали равна нулю. Поэтому движение брусков до и после касания пружины считаем равномерным. Для тел, составляющих замкнутую систему, можно применить закон сохранения энергии для колебаний пружины. При малой деформации пружины в ИСО колебания пружины с телами считаем гармоническими. 

Перейдем к решению. Для системы тел применим закон сохранения импульса при неупругом ударе: $m \vec v = ( m + M ) \vec u ,$ где u  — скорости тел после неупругого соударения. Учитывая, что $M=2m,$ получаем в проекциях на горизонтальную ось:

$m v =3mu,$ откуда $u= \frac{v}{3} .$

После соударения тела, как единое целое переместятся равномерно на расстояние L до касания с пружиной за время

\[t_1= \frac{L}{u} = \frac{3L}{v} .\]

Далее пружина сначала сожмется до максимальной деформации, а затем распрямится до недеформированного состояния. Время движения тел в этом случае равно $t_2= \frac{T}{2} ,$ где период колебаний пружинного маятника равен $T=2 \pi \sqrt{\frac{m + M}{k}} .$
Таким образом, время движения тел от момента касания до момента отрыва от пружины равно

\[t_2= \pi \sqrt{\frac{3m}{k}} .\]

При отсутствии силы трения по закону сохранения энергии скорость тел в момент отрыва от пружины равна скорости в момент касания u, причем отрыв произойдет в той же точке, где пружина снова станет недеформированной. Поэтому время движения тела вернутся в точку столкновения через промежуток времени t$_{1}$.
В итоге получаем общее время движения тел:

\[t=2t_1 + t_2= \frac{6L}{v} + \pi \sqrt{\frac{3m}{k}} = \frac{6 \cdot 0{,}15}{3} + 3{,}14 \cdot \sqrt{\frac{3 \cdot 0{,}1}{30}} \approx 0{,}614\,\text{с} .\]

Ответ: $t \approx 0{,}614\,\text{с}.$
}

\end{document}
